\documentclass[10pt,a4paper]{amsart}
\usepackage[margin=19mm]{geometry}
\usepackage{amsmath,amssymb,mathtools}
\usepackage[scr=rsfs,cal=euler]{mathalfa}
\usepackage{xfrac}
\usepackage[unicode,hidelinks,bookmarks=false]{hyperref}
\newcommand{\ie}{i.e.\ }
\newcommand{\cf}{cf.\ }
\newcommand{\resp}{respectively\ }
\newcommand{\Subsection}{\subsection}
\providecommand{\nobreakdash}{\nobreak\hbox{-}}
\setlength{\emergencystretch}{2em}
\multlinegap0pt
\usepackage{bm}
\usepackage{amssymb}
\usepackage{dsfont}
\newtheorem{lemma}{Lemma}
\title{Quantitative rapid stabilization for parabolic equations via the linear quadratic theory}
\author{Shengquan Xiang, Yu Xiao, Can Zhang}
\date{Source: arXiv:2606.09063v1 (CC BY 4.0)}
\begin{document}
\maketitle

\noindent\emph{Source and licence.} Quantitative rapid stabilization for parabolic equations via the linear quadratic theory. Version arXiv:2606.09063v1 (CC BY 4.0), distributed by arXiv under the Creative Commons Attribution 4.0 International licence, http://creativecommons.org/licenses/by/4.0/, as recorded in the arXiv licence metadata for this exact version and shown on the abstract page https://arxiv.org/abs/2606.09063v1. Full licence notice: \texttt{licenses/arxiv-2606.09063-CC-BY-4.0.txt}.\par\smallskip

\noindent\textit{Editorial scope.} Counted unit: the article's lemma PHIBOUND-NEW-18 (source0.tex lines 318--323). The article's complete original proof of it is delivered with it (source0.tex lines 325--363, one \texttt{proof} environment of 39 lines). No mathematics is written, repaired or re-derived: the statement and the proof are the article's own lines, in the article's own notation, and no retained line is substituted or re-flowed.  Retained uncounted as the source-local context the proof needs: lines 249--254; lines 266--268.  Nothing was omitted from any retained span, so the package carries no omission record.  No retained line contains a citation, so the wrapper carries no bibliography and no bibliography entry.  No retained line contains a figure environment, an image-inclusion command or drawing code, so no image is delivered and the package declares no asset dependency.  Editorial formatting only, in addition: an AMS \texttt{amsart} preamble that restates the article's own declarations and macros used by the retained lines, namely the environments \texttt{lemma} on separate counters, together with the standard packages those lines need.  The retained environments print in this document as Lemma 1 in order of appearance, whereas the article prints the same items with the numbering of its own full document.  The complete native source of the article as distributed in the arXiv e-print for this exact version is bundled unchanged as \texttt{sources/202378/source0.tex}.
\begin{equation}\label{TRANS-EQ}\begin{cases}
\partial_tz-\Delta z+az+b\cdot \nabla z-\lambda z=\chi_{\omega}v & \text{in } (0,+\infty)\times \Omega,\\[2mm]
z=0& \text{on }(0,+\infty)\times \partial\Omega,\\[2mm]
z( 0,\cdot)=y_0\in L^2(\Omega).
\end{cases}
\end{equation}

\begin{equation}\label{CO-FUNCAL-1}
\mathcal J_2(z,v;y_0)=\frac12\int_0^\infty\left(\|\nabla z(t)\|_{L^2(\Omega)}^2+\|v(t)\|_{L^2(\Omega)}^2\right)\,dt.
\end{equation} 

\begin{lemma}\label{PHIBOUND-NEW-18}
There exists a constant $\tilde{C}_2=\tilde{C}_2(\Omega,d,a,b)>0$ such that
\begin{equation}\label{PHI-bound-low-523} 
{\Phi_2} (y_0) \ge \tilde C_2\|y_0\|_{L^2(\Omega)}^2,\;\;\forall y_0\in L^2(\Omega).
\end{equation}
\end{lemma}

\begin{proof}
Since $\Phi_2(y_0)=\Phi_1(y_0)<\infty$, the standard minimizing sequence argument combined with the convexity of the cost functional \eqref{CO-FUNCAL-1} implies that $\Phi_2(y_0)$ is attained by a unique optimal pair $(z^*,v^*)$, so that the infimum is actually a minimum. 

Multiplying the equation \eqref{TRANS-EQ} by $2z^*$, we have
\begin{equation}\label{ENE-EST-1}
\begin{aligned}
\frac{d}{dt}\left(\|z^*(t)\|_{L^2(\Omega)}^2\right)&= -2\|\nabla z^*(t)\|_{L^2(\Omega)}^2
- 2\langle a z^*(t), z^*(t)\rangle_{L^2(\Omega)}- 2\langle b \cdot \nabla z^*(t), z^*(t)\rangle_{L^2(\Omega)}\\
&\quad+2\lambda\|z^*(t)\|_{L^2(\Omega)}^2+ 2\langle \chi_\omega v^*(t), z^*(t)\rangle_{L^2(\Omega)}.
\end{aligned}
\end{equation}
From
$$\frac{1}{2} \int_0^\infty  \left(\|\nabla z^*(t)\|_{L^2(\Omega)}^2 + \|v^*(t)\|_{L^2(\Omega)}^2\right) dt<\infty,$$
together with Poincar\'e's inequality, we conclude that
$\liminf\limits_{t\to \infty} \| z^*(t)\|_{L^2(\Omega)} =0$.
Hence, there exists an increasing sequence $\{T_n\}_{n\in\mathbb{N}^+}$ satisfying $T_n > T_{n-1}$ and $\|z^*(T_n)\|_{L^2(\Omega)} \le \frac{1}{n}, \forall\,n\in\mathbb{N}^+$.

Integrating \eqref{ENE-EST-1} over $(0,T_n)$, we obtain
\begin{equation}\label{ENE-EST-2}
\begin{aligned}
\|y_0\|_{L^2(\Omega)}^2 - \|z^*(T_n)\|_{L^2(\Omega)}^2&= 2\int_0^{T_n} \bigg(
\|\nabla z^*(t)\|_{L^2(\Omega)}^2+ \langle a z^*(t), z^*(t)\rangle_{L^2(\Omega)}+ \langle b\cdot\nabla z^*(t), z^*(t)\rangle_{L^2(\Omega)} \\
&\quad - \lambda \|z^*(t)\|_{L^2(\Omega)}^2- \langle \chi_\omega v^*(t), z^*(t)\rangle_{L^2(\Omega)}
\bigg)dt.
\end{aligned}
\end{equation}
Combined with $\|z^*(T_n)\|_{L^2(\Omega)} \le \frac{1}{n}$, this yields
\begin{equation*}
\|y_0\|_{L^2(\Omega)}^2\le \frac{1}{n^2}+ 2\int_0^{T_n} \bigg(\|\nabla z^*(t)\|_{L^2(\Omega)}^2+ \langle a z^*(t), z^*(t)\rangle_{L^2(\Omega)}+ \langle b\cdot\nabla z^*(t), z^*(t)\rangle_{L^2(\Omega)}- \langle \chi_\omega v^*(t), z^*(t)\rangle_{L^2(\Omega)}\bigg)dt.
\end{equation*}
Applying Young's inequality and taking $n\to\infty$, we deduce the lower bound
\begin{equation}\label{LOWER-BOUND}
\tilde C_2(\Omega,d,a,b)\|y_0\|_{L^2(\Omega)}^2
\le \frac12\int_0^\infty \big(
\|\nabla z^*(t)\|_{L^2(\Omega)}^2 + \|v^*(t)\|_{L^2(\Omega)}^2
\big)dt = \Phi_2(y_0).
\end{equation}
This completes the proof.
\end{proof}

\end{document}
